\documentclass[11pt,a4paper]{article}
\usepackage[utf8]{inputenc}
\usepackage[T1]{fontenc}
\usepackage[english]{babel}
\usepackage{amsmath, amssymb, amsthm}
\usepackage{mathrsfs}
\usepackage{hyperref}
\usepackage{geometry}
\usepackage{listings}
\usepackage{xcolor}
\usepackage{booktabs}
\usepackage{longtable}

\geometry{margin=2.5cm}

\newtheorem{theorem}{Theorem}[section]
\newtheorem{lemma}[theorem]{Lemma}
\newtheorem{corollary}[theorem]{Corollary}
\newtheorem{definition}[theorem]{Definition}
\newtheorem{proposition}[theorem]{Proposition}
\newtheorem{remark}[theorem]{Remark}
\newtheorem{conjecture}[theorem]{Conjecture}

\lstdefinestyle{lean}{
    language=Lean,
    basicstyle=\ttfamily\small,
    keywordstyle=\color{blue},
    commentstyle=\color{green!50!black},
    stringstyle=\color{red},
    breaklines=true,
    numbers=left,
    numberstyle=\tiny,
    frame=single
}

\title{Series XXXIII – Article XXXIII.0: Foundations of Spectral Proofs for Infinitely Many Prime Families}
\author{Christian Kithima \\
Independent Researcher, Kinshasa \\
\texttt{christiankithimaomba@gmail.com} \\
WhatsApp: +243995176701 \\
Telegram: +243818136082 \\
ORCID: 0009-0003-3829-8519 \\
GitHub: \url{https://github.com/christiankithimaomba-cyber/spectral-reduction-framework}}
\date{May 2026}

\begin{document}

\maketitle

\begin{abstract}
We extend the spectral method, already used to prove the infinitude of twin primes (Series VIII) and primes of the form \(n^2+1\) (Series XIV), to a large class of prime families defined by a fixed even gap \(d\). The key idea is the same finite verification (FV) strategy: prove an additive conjecture (every sufficiently large even integer is a sum of two primes that are both members of the family) using a sieve result (e.g., Chen's theorem in an arithmetic progression), then conclude the infinitude of the family. This article sets up the generic framework and outlines the series XXXIII.1–XXXIII.4. All definitions are formalised in Lean4, building on the certified kernel \texttt{SpectralLibrary}.
\end{abstract}

\tableofcontents

\section{Introduction}

The infinitude of twin primes \((p,p+2)\) was proved spectrally in Series VIII (Dubner). The method was then adapted to primes of the form \(n^2+1\) in Series XIV (Kithima‑Landau). In this series we unify and generalise the approach to any fixed even gap \(d\) for which an effective asymptotic bound exists (via Chen's sieve in an arithmetic progression). The result is a constructive proof that there are infinitely many prime pairs \((p,p+d)\).

\section{Recap of the spectral method for additive prime conjectures}

Let \(d\) be a fixed even integer. Define the property
\[
\text{prime\_pair}_d(p) \;:\; \text{Prime}(p) \land \text{Prime}(p+d).
\]
Consider the additive conjecture:
\[
\forall n \ge N_0,\; \exists p,q : \text{prime\_pair}_d(p) \land \text{prime\_pair}_d(q) \land p+q = n.
\]
Once this is proved (with an explicit \(N_0\)), the infinitude of prime pairs follows by the same argument as for twin primes.

The proof of the additive conjecture uses:
\begin{enumerate}
\item A boolean circuit \(C_{n,d}\) that tests, given binary representations of \(p,q\), whether both \(p\) and \(p+d\) are prime and \(q\) and \(q+d\) are prime and \(p+q=n\).
\item Reduction to 3‑SAT \(\Phi_{n,d}\) via Tseitin.
\item Construction of the spectral Hamiltonian \(H_{n,d} = \hat{V}_{n,d} + \Delta\).
\item Verification of the four pillars (P1–P4) – identical to the SAT case.
\item An effective asymptotic bound \(N_0(d)\) from analytic number theory (Chen's sieve in the appropriate arithmetic progression) ensuring that for \(n > N_0(d)\) the additive conjecture holds.
\item Finite verification for \(n \le N_0(d)\) via the D‑RSP algorithm.
\end{enumerate}

\section{Structure of the series XXXIII}

\begin{center}
\small
\begin{tabular}{@{}>{\bfseries}l>{\raggedright\arraybackslash}p{4.5cm}>{\raggedright\arraybackslash}p{6cm}@{}}
\toprule
\textbf{Article} & \textbf{Title} & \textbf{Main content} \\
\midrule
XXXIII.0 & Foundations of Spectral Proofs for Infinitely Many Prime Families & This article: introduction, plan, generic framework. \\
XXXIII.1 & Prime Cousins (gap 4): Boolean Circuit and Asymptotic Bound & Circuit for \(p+4\) primality, Chen's sieve in \(1\bmod4\), explicit \(N_0\), spectral verification. \\
XXXIII.2 & Sexy Primes (gap 6): Adaptation to Progression \(1\bmod6\) & Same for \(d=6\). \\
XXXIII.3 & Generalisation to Any Even Gap \(d\) & Conditions, parametrised circuit, uniform Lean proof. \\
XXXIII.4 & Synthesis – Infinitude of Several Prime Families & Summary table, integration into the global corpus of 33 problems. \\
\bottomrule
\end{tabular}
\normalsize
\end{center}

\section{Integration into the global corpus}

With this series, the list of 33 mathematical problems resolved by the spectral reduction lemma now includes entry \#33: Infinitude of prime cousins, sexy, and general even gaps.

\section{Formalisation in Lean4 (overview)}

The generic Lean code for any gap \(d\) will be derived from the twin prime code (Series VIII) by replacing the constant 2 with \(d\) and adapting the arithmetic progression for the sieve. The kernel provides the generic four pillars.

\bibliographystyle{plain}
\begin{thebibliography}{9}
\bibitem{kithimaVIII0}
  Kithima, C. (2026). Series VIII: Dubner's Conjecture and the Twin Prime Conjecture.
\bibitem{kithimaXIV0}
  Kithima, C. (2026). Series XIV: Kithima‑Landau Conjecture (Fourth Landau Problem).
\bibitem{chen}
  Chen, J. R. (1966). On the representation of a large even integer as the sum of a prime and the product of at most two primes.
\bibitem{lean}
  The Lean Community (2026). Lean 4 Theorem Prover Documentation.
\end{thebibliography}
\end{document}